\section{Postać Jordana macierzy.}

W tej sekcji nadal zakładamy, że pracujemy nad ciałem $\mathbb{C}$.

\begin{definition}
	Dla dowolnego $m \in \mathbb{Z}_+$ i dowolnego $\lambda \in \mathbb{C}$
		\textit{klatką Jordana} nazywamy macierz $J_m(\lambda) \in \mathbb{C}^{m \times m}$
	postaci
	\begin{equation}
		J_m(\lambda) = \left[\begin{matrix}
			\lambda & 1 & 0 & \dots & 0 \\
			0 & \lambda & 1 & \dots & \vdots \\
			0 & 0 & \lambda & \dots & 0 \\
			\vdots & \vdots & \vdots & \ddots & 1 \\
			0 & \dots & 0 & 0 & \lambda
		\end{matrix}\right].
	\end{equation}
\end{definition}

Okazuje się,
	że nawet jeśli macierz nie jest diagonalizowalna,
	to zawsze można ją sprowadzić do postaci,
	która jest prawie diagonalna.
Jednak nawet dla macierzy nad $\mathbb{R}$,
	ich postać Jordana może zawierać liczby zespolone.
Postać ta nazywa się postacią Jordana macierzy.
\begin{definition}
	Mówimy, że macierz $J$ jest w \textit{postaci Jordana},
		jeśli jest postaci
	\begin{equation}
		J = \left[\begin{matrix}
			J_{m_1}(\lambda_1) & \textbf{0} & \dots & \textbf{0} \\
			\textbf{0} & J_{m_2}(\lambda_2) & \dots & \textbf{0} \\
			\vdots & \vdots & \ddots & \vdots \\
			\textbf{0} & \textbf{0} & \dots & J_{m_r}(\lambda_r)
		\end{matrix}\right],
	\end{equation}
	gdzie każdy blok $J_{m_i}(\lambda_i)$ jest klatką Jordana.
\end{definition}
Postać normalna nie jest jednak unikalna!
Kolejność klatek w macierzy jest dowolna---zmienia to tylko permutację wymiarów w bazie.

Istnieje algorytm wyznaczania postaci normalnej Jordana.
Proces ten jest nieco bardziej złożony,
	niż proces wyznaczania macierzy diagonalnej.

% TODO

\begin{enumerate}
	\item oblicz wartości własne $A$: $\lambda_1, \dots, \lambda_k$,
	\item oblicz krotności algebraiczne $\lambda_1, \dots, \lambda_k$
		(krotności zer w wielomianie charakterystycznym $A$),
	\item oblicz krotności geometryczne $\lambda_1, \dots, \lambda_k$
		(krotność geometryczna $\lambda_i$ to $\dim(\ker(A - \lambda_i I)))$),
	\item (opcjonalnie) oblicz $\dim(\ker((A - \lambda_i I)^2)), \dim(\ker((A - \lambda_i I)^3)), \dim(\ker((A - \lambda_i I)^4)), \dots$, aż wymiar przestanie się zmieniać.
\end{enumerate}

\begin{example}
\begin{example}
	Wyznaczmy postać Jordana macierzy
	\begin{equation}
		A = \left[\begin{matrix}
			2 & 0 & 1 & 0 \\
			0 & 1 & 0 & 0 \\
			0 & 0 & 2 & 1 \\
			0 & 0 & 0 & 2
		\end{matrix}\right].
	\end{equation}
	Ponieważ $A$ jest macierzą górnotrójkątną,
		jej wartościami własnymi są $1$ i $2$.
	Ich krotności algebraiczne są równe odpowiednio $1$ i $3$.
	Aby ustalić postać Jordana dla wartości własnej $2$,
		obliczmy jej krotność geometryczną.
	Mamy
	\begin{equation}
		A - 2I = \left[\begin{matrix}
			0 & 0 & 1 & 0 \\
			0 & -1 & 0 & 0 \\
			0 & 0 & 0 & 1 \\
			0 & 0 & 0 & 0
		\end{matrix}\right],
	\end{equation}
	zatem
	\begin{equation}
		\ker(A - 2I) = \left\{\left[\begin{matrix}
			x_1 \\
			0 \\
			0 \\
			0
		\end{matrix}\right] : x_1 \in \mathbb{C}\right\}
		\implies
		\dim(\ker(A - 2I)) = 1.
	\end{equation}
	Krotność geometryczna wartości własnej $2$ jest więc równa $1$.
	Oznacza to,
		że dla wartości własnej $2$ występuje dokładnie jedna klatka Jordana.
	Ponieważ jej rozmiar musi być równy krotności algebraicznej wartości własnej $2$,
		ma ona rozmiar $3$.
	Dla wartości własnej $1$ mamy krotność algebraiczną $1$,
		więc odpowiada jej jedna klatka rozmiaru $1$.
	Postać Jordana macierzy $A$ jest więc równa
	\begin{equation}
		J = \left[\begin{matrix}
			2 & 1 & 0 & 0 \\
			0 & 2 & 1 & 0 \\
			0 & 0 & 2 & 0 \\
			0 & 0 & 0 & 1
		\end{matrix}\right].
	\end{equation}
\end{example}

\subsection{Zadania.}
